\clearpage
\section*{چکیده}

یادگیری فدرال یکی از رویکردهای نوین در حوزه یادگیری ماشین توزیع‌شده است که در آن، مدل هوش مصنوعی بدون نیاز به تبادل داده‌های خام، به‌صورت مشترک میان چندین گره محاسباتی آموزش می‌بیند. با گسترش روزافزون محیط‌های ابری چندگانه و ضرورت پردازش داده در لبه شبکه، پیاده‌سازی یادگیری فدرال در بسترهای ناهمگون از اهمیت ویژه‌ای برخوردار شده است. در چنین محیط‌هایی، گره‌های شرکت‌کننده از توان محاسباتی متفاوت، پهنای باند ارتباطی ناهمسان و توزیع داده غیریکنواخت برخوردارند که هر یک از این عوامل می‌توانند بر سرعت همگرایی و دقت مدل نهایی تأثیر منفی بگذارند.

در این پژوهش، یک مکانیزم اجماع سبک‌وزن تحت عنوان \lr{Trust-Weighted Adaptive Consensus (TWAC)} برای بهبود فرایند تجمیع در یادگیری فدرال پیشنهاد می‌شود. این مکانیزم با نگهداری امتیاز اعتماد پویا برای هر گره، به‌گونه‌ای عمل می‌کند که سهم گره‌های مخرب یا پراختلال در مدل نهایی کاهش یابد، بدون آنکه هزینه محاسباتی قابل‌توجهی به سیستم تحمیل شود. امتیاز اعتماد هر گره بر پایه دو مؤلفه به‌روزرسانی می‌شود: همسویی جهتی به‌روزرسانی آن گره با میانگین متحرک نمایی به‌روزرسانی‌های تجمیع‌شده گذشته، و منطقی‌بودن بزرگی به‌روزرسانی در مقایسه با میانه بزرگی به‌روزرسانی سایر گره‌ها. این ترکیب موجب می‌شود که روش پیشنهادی نسبت به هر دو نوع حمله جهت‌معکوس‌سازی و مقیاس‌افزایی به‌طور هم‌زمان مقاوم باشد.

برای ارزیابی روش پیشنهادی، یک بستر شبیه‌سازی محیط ابر چندگانه بر اساس مجموعه داده \lr{CIFAR-10} پیاده‌سازی شده است. در این بستر، بیست گره با توان محاسباتی و ویژگی‌های شبکه‌ای متفاوت شرکت دارند و توزیع داده میان آن‌ها از روش دیریکله با پارامتر $\alpha = 0.5$ پیروی می‌کند. نتایج آزمایش‌ها نشان می‌دهد که روش \lr{TWAC} در سناریوهای بدون حمله عملکردی معادل \lr{FedAvg} دارد، در حالی که در حضور گره‌های مخرب تا $30\%$، دقت بالاتری نسبت به روش‌های پایه به دست می‌آورد. همچنین، پیچیدگی محاسباتی این روش از مرتبه $O(nd)$ است که با \lr{FedAvg} برابر و از \lr{Krum} به‌مراتب کمتر است.

\vskip 5mm
\noindent\textbf{کلمات کلیدی}:
یادگیری فدرال، محیط ابر چندگانه، ناهمگونی محاسباتی، ناهمگونی شبکه، اجماع سبک‌وزن، مقاوم‌سازی در برابر گره‌های مخرب، توزیع داده غیریکنواخت.